\documentclass[a4paper,10pt]{article}
\usepackage{paper-en}
\usepackage{acts}
\usepackage{hyperref}

%\usepackage[notref,notcite,color]{showkeys}

\def\thetitle{Braid spaces have no geodesic loops}
\def\theauthors{Dmitri Panov and Anton Petrunin}

\hypersetup{colorlinks=true,
citecolor=black,
linkcolor=black,
anchorcolor=black,
filecolor=black,
menucolor=black,
urlcolor=black,
pdftitle={\thetitle},
pdfauthor={\theauthors}
}

%\overfullrule=100mm

\begin{document}

\title{\thetitle}
\author{\theauthors}
\date{}
\maketitle

\begin{abstract}
The $n$-string braid space is defined as a universal covering of the complex $n$-space branching at the hyperplanes $z_i=z_j$.
It is conjectured that the $n$-string braid space is $\CAT(0)$ for any $n$.
We describe a source of convex functions in the braid space and use them to prove that braid spaces have no geodesic loops.
\end{abstract}

\section{Introduction}

We assume that the reader is friends with $\CAT(0)$ spaces.

Let $\CC^n$ be the complex space
with coordinates $z_1,\ldots,z_n$ and the standard Euclidean metric.
Remove from $\CC^n$ the complex hyperplanes $z_i=z_j$ for all $i\ne j$;
denote the remaining set by $\WW^n$.
Let $\tilde \WW^n$ be the universal cover of $\WW^n$, equipped with the induced metric.
The completion of $\tilde \WW^n$ will be called the \emph{braid space}
with $n$ strings and denoted by $\CB^n$.

The set $\WW^n$ can be interpreted as the configuration space of $n$ punctures in the plane.
Therefore a path in $\WW^n$ describes a braid with $n$ strings, which explains our terminology.

The following conjecture is a special case of Daniel Allcock's conjecture \cite[Conjecture 1.4]{allcock}, which is related to an earlier conjecture of Ruth Charney and Michael Davis \cite[Conjecture 3]{charney-davis-95}, which in turn is motivated by the $K(\pi,1)$-conjecture \cite{vanderlek,paris}.
The authors are responsible for a more general version of Allcock's conjecture \cite[Conjecture 1.1]{panov-petrunin}, which we have been trying to prove for several years.

\begin{thm}{Conjecture}\label{braid-conj}
$\CB^n$ is $\CAT(0)$ for any integer $n\ge 1$.
\end{thm}

The case $n=2$ is easy; $\CB^2$ is isometric to a product of the Euclidean plane and a Euclidean cone over the real line.
The case $n=3$ is already nontrivial;
two proofs are given by the authors \cite{panov-petrunin};
one of them follows from the result of Ruth Charney and Michael Davis \cite{charney-davis-93}.
Ruth Charney also proved a special case of the conjecture for $n=4$ \cite[Theorem~5.4]{charney2004}.
Namely, she proved that the real braid space with 4 strings is $\CAT(0)$.
One may think of the real braid space as the space of braid diagrams.
Formally, it is the inverse image of $\RR^n$ under the natural projection $\CB^n\to \CC^n$.
This is a convex subset of $\CB^n$; so if the braid space is $\CAT(0)$, then so is the corresponding real braid space.

In this note we construct a class of convex functions on the braid space, which will be called \emph{pull-tight}.
The idea of this construction is inspired by Ian Agol's answer to our question \cite{agol}.

Let us describe pull-tight functions informally; an actual definition is given in Section~\ref{sec:pull-tight}.
Recall that $\WW^n$ is a configuration space of punctures in the plane.
Choose an isotopy class of a simple closed curve in the punctured plane and define a function $\tilde \WW^n\to\RR$ that returns the greatest lower bound of lengths of curves in this class;
this function extends naturally to $\CB^n$.
There is also a variant of this function for simple open curves that connect two fixed ideal points.

Here is the main result of our note.

\begin{thm}{Pull-tight theorem}\label{thm:pull-tight}
Every pull-tight function $\CB^n\to\RR$ is lower semicontinuous and convex.
\end{thm}

If $\CB^n$ is $\CAT(0)$, then it has no geodesic digons.
(Since $\CB^n$ is a polyhedral space and it has finitely many shapes, the converse holds as well; see \cite[II.5.4]{bridson-haefliger} and also compare with \cite[3.4.3]{AKP-invitation}.)
The pull-tight theorem implies the following weaker statement.

\begin{thm}{Theorem}\label{thm:no-loops}
The braid space has no geodesic loops.
\end{thm}

In the proof we assume that a geodesic loop exists, and construct a pull-tight function that is nonnegative, not constant on the loop, and vanishes at the base point; this contradicts the pull-tight theorem.

\section{Number of collisions}

Consider a geodesic $\gamma$ in $\CB^n$.
A point $\gamma(t)$ defines an array of punctures $p_1(t),\z\ldots,p_n(t)\in\CC$.
We say that $\gamma$ does not have a  \emph{collision} at time $t$ if the velocities of all $p_i$ are constant in a neighborhood of $t$,
so the punctures move with constant velocity between the collision times.
(A priori, the set of collision times is merely closed.)

Since $\CB^n$ is a polyhedral space with finitely many shapes, the following lemma follows from the theorem of Martin R. Bridson and André Haefliger \cite[I.7.28]{bridson-haefliger}.

\begin{thm}{Lemma}\label{lem:BFK}
There is a universal bound $N=N(n)$ on the number of collisions of geodesics in $\CB^n$.
\end{thm}


Let us mention that this lemma is also a limiting case of the result of Dmitri Burago, Serge Ferleger and Alexey Kononenko on the number of collisions in semi-dispersing billiards \cite{burago-ferleger-kononenko-1997,AKP-invitation};
their proof adapts to our situation and gets simpler.


\section{Pull-tight functions}\label{sec:pull-tight}

A \emph{pull-tight} function is a function $f\colon\CB^n\to\RR$ that can be obtained via the following procedure.
We first define it on the universal covering $\tilde\WW^n$ of the space $\WW^n$ and then extend it to $\CB^n$.

\begin{wrapfigure}{o}{40mm}
\centering
\vskip-0mm
\includegraphics{mppics/pic-70}
\vskip-0mm
\end{wrapfigure}

A point $\bm{p}\in\WW^n$ will be regarded as an array of $n$ distinct punctures $(p_1,\ldots,p_n)$ of the complex plane $\CC$.
Choose a simple oriented curve $\lambda$ in the punctured plane.
The curve can be closed or open; in the latter case, we assume that $\lambda$ connects two fixed ideal points.

Think of the curve as a thread and make it tight.
The curve becomes a polygonal line, say $\ell$, with vertices at the punctures; self-intersections are possible.
If $\lambda$ is open, then $\ell$ has two infinite edges.
The polygonal line $\ell$ can be arbitrarily well approximated by a simple curve isotopic to~$\lambda$.

If $\lambda$ is closed, then $f(\bm{p})$ is defined as the length of $\ell$.

If $\lambda$ is open, then $\ell$ might not exist.
It happens if $\lambda$ can be pushed into a strictly convex set containing no punctures, so it is possible to shrink the thread indefinitely.
In this case, we assume that the pull-tight function is undefined (one could also define it as $f(\bm{p})\equiv -\infty$).
Otherwise $f(\bm{p})$ is defined as the excess of its length compared to the two-edge curve passing thru the origin.
More formally, suppose that $\alpha$ and $\omega$ are the ideal points at the beginning and the end,
and $p_{i_1},\ldots,p_{i_k}$ is the sequence of vertices of $\ell$.
Then
\[f(\bm{p})
=
|\alpha-p_{i_1}|^\flat+|p_{i_1}-p_{i_2}|+\ldots+|p_{i_{k-1}}-p_{i_k}|+|p_{i_k}-\omega|^\flat,\]
where $|\xi-x|^\flat$
denotes the Busemann function at $x$ with ideal point $\xi$ normalized so that $|\xi-0|^\flat=0$; so,
\[|\xi-x|^\flat=-\langle u_\xi,x\rangle,\]
where $u_\xi$ denotes the unit vector in the direction of $\xi$.
Note that $|\xi-x|^\flat$ might take negative values.

In both cases we have defined a function on $\tilde \WW^n$.
Note that it is locally Lipschitz, and it naturally extends to a lower semicontinuous function on $\CB^n$.
Along a geodesic this extension is continuous: at a collision the edges inside the colliding clusters shrink to zero, while all remaining terms have their usual limits.

\parit{Proof of \ref{thm:pull-tight}.}
Let $f\:\CB^n\to \RR$ be a pull-tight function, and let $\lambda$ be its defining curve.
Given a unit-speed geodesic $\gamma$ in $\CB^n$,
we need to show that $f\circ\gamma$ is a convex real-to-real function.

Each point $\gamma(t)$ defines an array of punctures $p_1(t),\ldots,p_n(t)\in\CC$.
The curve $\lambda$ divides the punctures into two sets that lie on its left and right sides.

Denote by $\ell(t)$ the pull-tight curve at time $t$.
It is a polygonal line that can be encoded by a sequence of its vertices; each vertex must be one of $p_i(t)$;
self-intersections, repeated vertices and edges are allowed.
By construction, there is a continuous one-parameter family of arbitrarily small perturbations of $\ell(t)$ and $p_i(t)$ to simple curves isotopic to $\lambda$ in the punctured plane with $n$ distinct punctures.


Suppose that the vertex sequence of $\ell(t)$ does not change in some open time interval $\II$ and no collision occurs in $\II$.
Note that if two points move with constant velocities in the plane, then the distance between them is a convex function of time.
It follows that $f\circ\gamma|_{\II}$ is convex.

The vertex sequence can change in two ways:
\begin{enumerate}[(i)]
 \item when a puncture hits or leaves $\ell(t)$
and
 \item at a collision.
\end{enumerate}
Such events happen at finitely many times, say $t_1,\ldots,t_m$.
Indeed, by \ref{lem:BFK}, we have only finitely many collisions,
and between the collision times hitting/leaving events are described by isolated zeros of algebraic equations.

For each $t_i$, we have to check that
\[(f\circ\gamma)'(t_i+) - (f\circ\gamma)'(t_i-)\ge 0;
\eqlbl{1star}\]
here $(f\circ\gamma)'(t\pm)$ denote one-sided derivatives.
Once this is done, convexity of $f\circ\gamma$ follows.


\parit{Case (i).}
Suppose that the vertex sequence changes at time $t_i$ without a collision.
For example, $\ell(t)$ catches up with the point $p_5(t)$, so $p_5(t)$ enters the vertex sequence between vertices $p_2(t)$ and $p_7(t)$, as in the picture.

\begin{wrapfigure}{o}{46mm}
\centering
\vskip-0mm
\includegraphics{mppics/pic-50}
\vskip-0mm
\end{wrapfigure}

In a neighborhood of $t_i$ the curve is described by two vertex sequences;
denote the corresponding functions by $f_{\textrm{before}}$ and $f_{\textrm{after}}$;
if $t\approx t_i$, then
\[f\circ\gamma(t)=
\begin{cases}
f_{\textrm{before}}\circ\gamma(t)&\text{if\ } t\le t_i,
\\
f_{\textrm{after}}\circ\gamma(t)&\text{if\ } t\ge t_i.
\end{cases}
\]
Note that
\[(f_{\textrm{before}}\circ\gamma)'(t_i)=(f_{\textrm{after}}\circ\gamma)'(t_i);\]
hence \ref{1star} follows.

\parit{Case (ii).} Suppose that $k$ punctures collide at time $t_i$.
We can assume that $t_i=0$ and the collision point is also $0\in\CC$.

\begin{figure}[b!]
\centering
\vskip-0mm
\includegraphics{mppics/pic-60}
\vskip-0mm
\end{figure}

Rescale a neighborhood of the collision point and a time interval around the collision time, keeping the velocities the same.
In the limit we get a collision of all $k$ punctures in $\CB^k$.
The punctures come with constant velocities to $0$ and run from it with constant velocities.
Denote by $\tilde \ell(t)$, $\tilde \gamma$, and $\tilde p_1,\ldots,\tilde p_k$
the limits of the pull-tight curve, the geodesic in $\CB^k$, and the corresponding punctures.
Suppose that $\tilde \ell(t)$ is still defined by the isotopy class of one simple curve $\tilde\lambda$ in the $k$-punctured plane with fixed ideal endpoints.
Let us check \ref{1star} for the corresponding pull-tight function $\tilde f\colon\CB^k\to \RR$.

Note that the shape of $\tilde \ell(t)$ changes only at $t=0$:
for $t>0$ and $t<0$ it only rescales.
At $t=0$ these shapes approach a plane angle from both sides;
let $\theta$ be its angle measure.
If \ref{1star} does not hold, then $\theta<\pi$; in other words, the ideal points of $\tilde \ell$ cannot be opposite.
In this case, adding a constant vector to all velocities in the $k$-puncture array can reduce the problem to the case
\[(\tilde f\circ\tilde \gamma)'(0-)>0>(\tilde f\circ\tilde \gamma)'(0+).\]
We may still assume that the collision happens at zero at time zero;
so, $\tilde f\circ\tilde \gamma(0)\z=0$ and from above
\[\tilde f\circ\tilde \gamma(t)\le -\eps\cdot |t|\eqlbl{sharp-max}\]
for some $\eps>0$.
Now we are going to construct a shortcut of $\tilde \gamma$, showing that it is not geodesic.

The curve $\tilde\ell(t)$ might have self-intersections, but it admits an approximation by a simple curve, which divides the plane into two regions on its left and right sides. %Чуть выше мы уже сказали, что у неё только одна компонента.
Let us equip these regions with the induced length metric;
they are $\CAT(0)$ \cite[7.12]{AKP-invitation}.
Pass to their limits as the curve approaches $\tilde\ell(t)$; denote the limit spaces by $L_t$ and $R_t$; they are $\CAT(0)$ as limits of $\CAT(0)$ spaces.%
\footnote{If $\tilde\ell(t)$ is simple, then it bounds $L_t$ and $R_t$; our construction generalizes this picture.}
One of the angles $L_0$ or $R_0$ is a convex angle of measure $\theta$;
let us assume that this is~$L_0$.

Choose an ideal point $\xi$ that lies midway between the ideal points, say $\alpha$ and $\omega$, of $\tilde\ell(t)$.
Let us denote by $|x-\xi|^\flat_t$ its Busemann function $L_t\to\RR$.
These functions are defined in the length metric of $L_t$.
Formally speaking, $\xi$ corresponds to an ideal point in each space $L_t$,
but far from $0$, the half-lines in $L_t$ in the direction of $\xi$ become parallel.
Thus the Busemann functions $x\z\mapsto |x-\xi|^\flat_t$ can be normalized uniformly.
(Far away from $0$, they all agree with the Busemann function for $\xi$ with respect to the Euclidean metric: $x\mapsto-\langle x,u_\xi\rangle$.)
In particular we have $|t\cdot x-\xi|^\flat_t=t\cdot |x-\xi|^\flat_1$ for any $t\ge 0$ and $x\in L_1$.
The superlevel sets of this Busemann function are compact sets in $L_t$.
Moreover,
\[2\cdot \cos \tfrac\theta2 \cdot |x-\xi|^\flat_t
\le
|\alpha-x|^\flat_t+|x-\omega|^\flat_t
\le \tilde f\circ \tilde\gamma(t)\]
for any $x\in L_t$;
the first inequality is $\CAT(0)$ comparison and the second is the triangle inequality, both applied to far points in directions of $\xi$, $\alpha$, and $\omega$.

Since $\tilde f\circ\tilde \gamma(t)\le -\eps\cdot |t|$, we get
\[|x-\xi|^\flat_t\le -\eps'\cdot|t|\]
for any $x\in L_t$ and some fixed $\eps'>0$.

Suppose $\tilde p_i$ lies on the left side of~$\tilde \lambda$.
Let $\delta_i=-|\tilde p_i(1)-\xi|^\flat_1$;
note that $\tilde p_i(t)\z=t\cdot \tilde p_i(1)$ and $|\tilde p_i(t)-\xi|^\flat_t=-t\cdot \delta_i$ for any $t\ge 0$.
Let us connect $\tilde p_i(1)$ to $\xi$ by a unit-speed geodesic half-line $\sigma_i\:[\delta_i,\infty)\z\to L_1$;
so $|\sigma_i(s)-\xi|^\flat_1+s\equiv0$.
Note that
\[s\mapsto t\cdot \sigma_i(s/t)
\eqlbl{2star}\]
defines the unit-speed geodesic half-line in $L_t$ from $\tilde p_i(t)$ to $\xi$.
Further, if $s$ is large, then
\[\sigma_i(s)=s\cdot u_\xi+ z_i,\eqlbl{3star}\]
for a complex constant $z_i$ such that $z_i\perp u_\xi$.

Choose $\delta>0$ such that $\delta<\delta_i$ for any $\tilde p_i$ on the left side of $\tilde \lambda$.
Let $\hat p_i(t)$ be $\tilde p_i(t)$ moved toward $\xi$ along the half-line defined by \ref{2star}:
if $\tilde p_i(t)$ is below the $(-\delta)$-level of the Busemann function for $\xi$, then $\tilde p_i(t)$ is not moved;
otherwise it is moved to the $(-\delta)$-level.
In other words, \[\hat p_i(t)=t\cdot \sigma_i(\max\{\,\delta/t, \delta_i\,\});\]
so, if $t\cdot \delta_i\le \delta$, then
we move $\tilde p_i(1)$ along $\sigma_i$ to the $(-\delta/t)$-level and then scale it down by $t$.

Note that $|\tilde p_i'(t)|^2=|\tilde p_i(1)|^2$ for any $t>0$.
Therefore
\[\int\limits_{0}^{\delta/\delta_i}|\tilde p_i'(t)|^2\cdot dt
=
(\delta/\delta_i)\cdot|\tilde p_i(1)|^2.\]
If $0<t<\delta/\delta_i$, then $\hat p_i(t)=t\cdot \sigma_i(\delta/t)$
and
\[|\sigma'_i|=1
\quad\Rightarrow\quad
|\hat p_i'(t)|^2=\left(\frac{|\hat p_i(t)|^2-\delta^2}{t}\right)'\]
almost everywhere.
By \ref{3star}, we have $|\hat p_i(t)|^2-\delta^2=o(t)$, whence
\[\int\limits_0^{\delta/\delta_i}|\hat p_i'(t)|^2\cdot dt
=
(\delta/\delta_i)\cdot(|\tilde p_i(1)|^2-\delta_i^2)
<
\int\limits_0^{\delta/\delta_i}|\tilde p_i'(t)|^2\cdot dt.\]

Observe that all modified punctures coincide at time $0$: if $\tilde p_i$ lies on the left side of $\tilde\lambda$, then $\hat p_i(0)=\delta\cdot u_\xi$.
Therefore, repeating the same construction for $t\le 0$ produces a curve $t\mapsto \hat p_i(t)$ defined on the whole $\RR$.

Each puncture $\tilde p_i$ on the right side will not be modified;
that is, we let $\hat p_i(t)=\tilde p_i(t)$.
Thus we obtain a new curve $t\mapsto \hat p_i(t)$ for each $i$.
These curves correspond to a new curve $\hat\gamma$ in $\CB^k$ that coincides with $\tilde \gamma$ away from some time interval $[-R,R]$.
Indeed, punctures cannot cross $\tilde\ell(t)$; so the punctures on its different  sides do not interact.
On the right side the punctures move as before,
and all modified punctures approach $\delta\cdot u_\xi$ at time $0$,
so we are free to jump to any branch at this moment.

The above inequality implies that
\[\int\limits_{-R}^{R}|\tilde p'_i(t)|^2\cdot dt
\ge
\int\limits_{-R}^{R}|\hat p'_i(t)|^2\cdot dt\]
for each $i$; moreover, if $\tilde p_i$ is on the left side then the inequality is strict.
Since $\tilde\gamma$ has constant speed,
summing over $i$ and applying Cauchy--Bunyakovsky, we obtain
\[\length \tilde\gamma|_{[-R,R]}>\length\hat\gamma|_{[-R,R]},\]
which gives the required contradiction.

The limit curve $\tilde\ell(t)$ might be empty, but then the collision makes no contribution to the derivative jump of $f\circ\gamma$.
Further, $\tilde\ell(t)$ might be closed, but then $\tilde f\circ \tilde\gamma$ cannot satisfy \ref{sharp-max}.
Furthermore, $\tilde\ell(t)$ might have several components;
they have to be described by isotopy classes of nonintersecting simple curves $\tilde\lambda_1,\ldots,\tilde\lambda_m$ with fixed ideal endpoints.
Denote by $\tilde f_1,\dots,\tilde f_m$ the corresponding pull-tight functions.
One of the functions $\tilde f_i\circ\tilde\gamma$ must have a negative derivative jump at $0$ (the original derivative jump is the sum of the component jumps).
Applying the common-velocity adjustment, we can assume that $\tilde f_i\circ\tilde\gamma$
has  a sharp maximum at zero; in other words, $\tilde f_i\circ\tilde\gamma$ satisfies \ref{sharp-max}.
Then we can repeat the argument for $\tilde\lambda_i$.
As before, we get a shortcut $\hat\gamma$ for $\tilde \gamma$ and arrive at a contradiction.

\parit{Mixed case.} Our calculations are local: combining them, one can prove \ref{1star} in the case when several events (collisions and hitting/leaving) happen simultaneously.
\qeds

\section{No geodesic loops}

\parit{Proof of \ref{thm:no-loops}.}
Assume $\gamma\:[0,L]\to\CB^n$ is a geodesic loop;
denote by $p_1(t),\z\ldots,p_n(t)\in\CC$ its punctures at the time $t$.

Choose $n-1$ pull-tight functions $f_1,\ldots,f_{n-1}$
such that the corresponding pull-tight curves are vertical straight lines that separate the punctures at time~$0$.

With the natural normalization we have
\[f_i\circ\gamma(0)=f_i\circ\gamma(L)=0\]
for each $i$.
Since $f_i$ are convex and nonnegative, $f_i\circ\gamma(t)\equiv 0$ for all $i$.
It follows that pull-tight lines form moving vertical barriers between the punctures giving them no chance to interact;
in particular, each puncture moves with constant velocity,
and the image of $\gamma$ in $\CC^n$ is a line segment --- a contradiction.
\qeds

\section{To whom it may concern}

The following observations might save time for someone interested in the conjecture; this section explains where we hit the wall.
Consider this part as a letter (while writing it we had in mind Ruth Charney and Michael Davis).

In the definition of $\CB^n$ one can lift an arbitrary diagonal metric on $\CC^n$.
In other words, we prescribe positive masses $\mu_1,\z\ldots,\mu_n$ to the punctures $p_1(t),\z\ldots, p_n(t)$ that move in $\CC$ and define the speed of the corresponding path in $\CB^n$ as
\[\sqrt{\mu_1\cdot|p_1'(t)|^2+\ldots+ \mu_n\cdot|p_n'(t)|^2}.\]
Let us denote the obtained space by $\CB^n(\bm{\mu})$, where $\bm{\mu}=(\mu_1,\ldots,\mu_n)$.%
\footnote{Equivalently, we can apply a diagonal matrix with $\sqrt{\mu_1},\ldots,\sqrt{\mu_n}$ on the diagonal to the complex hyperplane arrangement $z_i=z_j$ (so the new complex arrangement is defined by $\frac{z_i}{\sqrt{\mu_i}}=\frac{z_j}{\sqrt{\mu_j}}$).
Then $\CB^n(\bm{\mu})$ is isometric to the universal cover of $\CC^n$ (with the standard metric) infinitely branching along the new hyperplane arrangement.}

The definition of a pull-tight function extends to $\CB^n(\bm{\mu})$.
Moreover, the proof of the pull-tight theorem has a straightforward generalization that gives the following.

\begin{thm}{Observation}
Every pull-tight function $\CB^n(\bm{\mu})\to\RR$ is lower semicontinuous and convex.
\end{thm}

\begin{thm}{Corollary from the conjecture}\label{cor:masses}
If the conjecture holds, then $\CB^n(\bm{\mu})$ is $\CAT(0)$ for any array of positive masses $\bm{\mu}=(\mu_1,\ldots,\mu_n)$.
\end{thm}

\parit{Proof.}
If the masses $\mu_i$ are integers, then instead of the $i$-th string of the braid one can consider $\mu_i$ identical copies of this string.
If $M=\mu_1+\ldots+\mu_n$, then $\CB^n(\bm{\mu})$ is isometric to a convex subset of $\CB^M$.
In particular, the conjecture for $\CB^M$ implies the corollary for $\CB^n(\bm{\mu})$.

By rescaling, rational approximation, and passing to the limit we get the case of arbitrary positive masses.
\qeds

\parit{Approach to the conjecture.} Let us try to prove \ref{cor:masses} by induction on $n$.
The case $n=3$ already follows from \cite{charney-davis-93} (checking the angle condition is straightforward).
The first new case $n=4$ might go as follows.
Assuming the contrary, we have a braid space $\CB^4(\bm{\mu})$ with a geodesic bigon.

We can split off $\CC$ from $\CB^4(\bm{\mu})$, so that the center of mass of the punctures is at $0$.
Denote by $\Sigma^4(\bm{\mu})$ the unit sphere in the remaining factor.
We need to show that it is $\CAT(1)$.

By the induction hypothesis $\Sigma^4(\bm{\mu})$ is locally $\CAT(1)$.%
\footnote{In fact our theorem implies that every $\tfrac\pi2$-ball in $\Sigma^4(\bm{\mu})$ equipped with the induced length metric is $\CAT(1)$;
in particular, $\Sigma^4(\bm{\mu})$ is $\CAT(4)$ and any $\tfrac\pi4$-ball in $\Sigma^4(\bm{\mu})$ is $\CAT(1)$.}
Assume $\Sigma^4(\bm{\mu})$ is not $\CAT(1)$.
By \cite[II.5.4(7)]{bridson-haefliger}, it contains a shortest geodesic circle of length $L(\bm{\mu})<2\cdot\pi$ (a \emph{geodesic circle} is a distance-preserving embedding of a circle with length metric).
If one of the masses approaches zero, we expect to have $L(\bm{\mu})\z\to2\cdot \pi$ by the induction hypothesis.

Choose a geodesic circle of length $L(\bm{\mu})$ in $\Sigma^4(\bm{\mu})$.
This geodesic is not null-homotopic in the class of curves of length $<2\cdot\pi$.
This property survives after perturbing the metric.
The new curve might fail to be a geodesic, but its length gives an upper bound on $L(\bm{\mu})$ after small perturbation of masses.

Given an array of punctures $p_1,\ldots, p_4$, we translate it to make its center of mass zero, then normalize it so that $\mu_1\cdot|p_1|^2+\ldots+ \mu_4\cdot|p_4|^2=1$; it corresponds to a point in $\Sigma^4(\bm{\mu})$.
This construction identifies the underlying sets of $\Sigma^4(\bm{\mu})$ for different masses; so we may think of $\Sigma^4(\bm{\mu})$ as the same set equipped with different metrics.

Since the rescaling of masses does not change the isometry type of $\Sigma^4(\bm{\mu})$, we get
\[\mu_1\cdot \frac{\partial L}{\partial \mu_1}+\ldots+\mu_4\cdot \frac{\partial L}{\partial \mu_4}=0.\]

The function $\bm{\mu}\mapsto L(\bm{\mu})$ does not have to be smooth,
but it is semiconcave.
At a minimum point of $L$ (the induction hypothesis should imply that it is in the interior), we have $\frac{\partial L}{\partial \mu_i}=0$ for every $i$.
Assume that a unit-speed parametrization of the geodesic circle is given by punctures $p_1(t),\ldots,p_4(t)$; so
\begin{align*}
\mu_1\cdot p_1(t)+&\ldots+ \mu_4\cdot p_4(t)=0,
\\ \mu_1\cdot|p_1(t)|^2+&\ldots+ \mu_4\cdot|p_4(t)|^2=1,
\\ \mu_1\cdot|p_1'(t)|^2+&\ldots+ \mu_4\cdot|p_4'(t)|^2=1.
\end{align*}
Moreover, $p_i''(t)+p_i(t)=0$ if $p_i$ does not collide at time $t$.
The condition $\frac{\partial L}{\partial \mu_i}=0$ means that we have a zero sum of $\langle\, p_i(t_j),p_i'(t_j+)-p_i'(t_j-)\, \rangle$ at the collision times~$t_j$.

The last identity might contradict $L<2\cdot\pi$.
Also, the estimate on $L$, which is described above, might imply that (in the barrier sense) the Hessian of $L$ is negative in some direction, which should prove the conjecture.
AI suggests looking at the Laplacian with respect to the Fisher metric on the mass simplex: $g(h,h)=\tfrac{h_1^2}{\mu_1}+\ldots+\tfrac{h_4^2}{\mu_4}$ where $h_1+\ldots+h_4=0$, which makes it closely related to the description of $\CB^4(\bm{\mu})$ in the footnote.

Note however that for an equator the Hessian and Laplacian vanish.
Therefore, one has to use that $L<2\cdot\pi$ and the geometry of the geodesic circle.
The following observation might help: since $L<2\cdot\pi$, the image of the geodesic circle in $\mathbb{S}^5\subset \CC^3$ lies in a hemisphere.

\parbf{Lasso functions.}
The proof of convexity of pull-tight functions does not use much and has variations;
here is one example.

Choose one puncture $p_i$, called the \emph{caught puncture}.
Pass to the covering of $\CC$ infinitely branching at all punctures except $p_i$, and lift $p_i$ to this covering.
Consider the Busemann function from the chosen ideal point $\xi$ to the lift.
Our proof of convexity works for this function as well.
The obtained function $s$ (let's call it a \emph{lasso function})
has a closely related pull-tight function $f$ --- the one with both ideal endpoints at $\xi$ and $\lambda$ surrounding just $p_i$.

Note that $2\cdot s\le f$; in particular, the subzero set of $s$ contains that of $f$.
So, in a sense $s$ looks more like a Busemann function on $\CB^n$.
(If Busemann functions could be approximated by convex functions, then the conjecture would follow, but one may need much less.)

{\sloppy
\def\emph{\textit}
\printbibliography[heading=bibintoc]
\fussy
}

\end{document}


